\documentclass[12pt,a4paper]{amsart}

\usepackage{amsmath,amssymb,amsthm}
\usepackage{mathtools}
\usepackage{hyperref}
\usepackage{geometry}
\geometry{margin=2.5cm}

% -------------------------------------------------------
% Theorem environments
% -------------------------------------------------------
\newtheorem{theorem}{Theorem}[section]
\newtheorem{lemma}[theorem]{Lemma}
\newtheorem{corollary}[theorem]{Corollary}
\newtheorem{proposition}[theorem]{Proposition}
\theoremstyle{definition}
\newtheorem{definition}[theorem]{Definition}
\newtheorem{remark}[theorem]{Remark}

ewtheorem{example}[theorem]{Example}

% -------------------------------------------------------
% Custom commands
% -------------------------------------------------------
\newcommand{\N}{\mathbb{N}}
\newcommand{\Z}{\mathbb{Z}}
\newcommand{\Q}{\mathbb{Q}}
\newcommand{\Fp}{\mathbb{F}_p}
\DeclareMathOperator{\ord}{ord}
\newcommand{\Wq}[1]{W_{#1}}

% -------------------------------------------------------
\begin{document}
% -------------------------------------------------------

\title{The Wilson Quotient and Wilson Primes}

\author{Michael Fuhrmann}
\address{Specialist Mathematics Project, Build 64}
\email{specialist-maths@localhost}
\date{\today}

\begin{abstract}
Wilson's theorem guarantees that $(p-1)! + 1$ is divisible by every prime $p$.
This motivates the \emph{Wilson quotient} $W_p = \bigl((p-1)!+1\bigr)/p \in \Z$.
A prime $p$ is called a \emph{Wilson prime} if $p^2 \mid (p-1)!+1$, equivalently
if $p \mid W_p$.  We prove four main results:
(i) an explicit congruence expressing $W_p$ in terms of the harmonic sum
$\sum_{j=1}^{p-1} j^{-1} \pmod p$,
(ii) the equivalence of the Wilson prime condition with $p^2 \mid (p-1)!+1$,
(iii) a connection to Bernoulli numbers via Kummer's congruence
$W_p \equiv -B_{p-1} \pmod p$,
and (iv) a heuristic argument for the rarity of Wilson primes.
The only known Wilson primes are $5$, $13$, and $563$; no further Wilson prime
exists below $2 \times 10^{13}$.
\end{abstract}

\maketitle
\tableofcontents

% ================================================================
\section{Introduction}
% ================================================================

By Wilson's theorem, $(p-1)! \equiv -1 \pmod p$ for every prime $p$.
This means that $p \mid (p-1)! + 1$, so the integer
\[
  W_p = \frac{(p-1)! + 1}{p}
\]
is well-defined.  The question of whether $p$ further divides $W_p$ is the
starting point of the theory of Wilson primes.

Wilson primes are analogous to Wieferich primes in Fermat's little theorem:
just as Wieferich primes satisfy $a^{p-1} \equiv 1 \pmod{p^2}$ (the
``squared'' version of Fermat), Wilson primes satisfy the ``squared'' version
of Wilson.  Both are extremely rare, and in both cases only finitely many are known.

Throughout, $p$ denotes a prime and all arithmetic in residues is taken
modulo the prime written after $\pmod{}$.

% ================================================================
\section{The Wilson Quotient and Its Basic Properties}
% ================================================================

\begin{definition}[Wilson quotient and Wilson primes]\label{def:wilson_quotient}
  For a prime $p$, the \emph{Wilson quotient} is
  \[
    W_p \;=\; \frac{(p-1)!\,+\,1}{p} \;\in\; \Z.
  \]
  The prime $p$ is a \emph{Wilson prime} if $p \mid W_p$,
  equivalently if $p^2 \mid (p-1)! + 1$.
\end{definition}

\begin{proposition}[Integrality and basic values]\label{prop:basic}
  $W_p$ is a positive integer.  The first few values are:
  \[
    W_2 = 1,\quad W_3 = 1,\quad W_5 = 5,\quad W_7 = 103,\quad W_{11} = 329891.
  \]
\end{proposition}

\begin{proof}
  By Wilson's theorem, $p \mid (p-1)! + 1$, so $W_p \in \Z$.
  Positivity follows from $(p-1)! + 1 \geq 1$.
  The numerical values are direct computation:
  \begin{align*}
    W_2 &= \tfrac{1!+1}{2} = 1,\\
    W_3 &= \tfrac{2!+1}{3} = 1,\\
    W_5 &= \tfrac{4!+1}{5} = \tfrac{25}{5} = 5,\\
    W_7 &= \tfrac{6!+1}{7} = \tfrac{721}{7} = 103,\\
    W_{11} &= \tfrac{10!+1}{11} = \tfrac{3628801}{11} = 329891.\qedhere
  \end{align*}
\end{proof}

\begin{theorem}[Wilson prime criterion]\label{thm:wilson_prime_criterion}
  A prime $p$ is a Wilson prime if and only if $W_p \equiv 0 \pmod p$,
  i.e.\ $p^2 \mid (p-1)! + 1$.
\end{theorem}

\begin{proof}
  By definition, $W_p = \bigl((p-1)!+1\bigr)/p$.  Then
  \[
    p \mid W_p
    \iff p \mid \frac{(p-1)!+1}{p}
    \iff p^2 \mid (p-1)! + 1. \qedhere
  \]
\end{proof}

% ================================================================
\section{A Half-Factorial Formula for the Wilson Quotient}
% ================================================================

The following theorem gives a concrete formula for $W_p \pmod p$ in terms
of the half-factorial $m! = ((p-1)/2)!$ and a half-harmonic sum.

\begin{theorem}[Wilson quotient via pairing]\label{thm:harmonic}
  Let $p$ be an odd prime and set $m = (p-1)/2$.  Define
  \[
    \mu_p \;=\; \frac{(-1)^m (m!)^2 + 1}{p} \;\in\; \Z
    \qquad\text{and}\qquad
    S_p \;=\; \sum_{j=1}^{m} j^{-1} \pmod p,
  \]
  where $j^{-1}$ is the multiplicative inverse of $j$ modulo $p$.  Then
  \[
    W_p \;\equiv\; \mu_p + S_p \pmod p.
  \]
  In particular, $p$ is a Wilson prime if and only if $\mu_p + S_p \equiv 0 \pmod p$.
\end{theorem}

\begin{remark}
  Note carefully that the full sum $\sum_{j=1}^{p-1} j^{-1} \equiv 0 \pmod p$
  always holds (since $\{j^{-1}\}_{j=1}^{p-1} = \{1,\ldots,p-1\}$ as sets mod $p$),
  so only the \emph{half}-sum $S_p = \sum_{j=1}^{(p-1)/2} j^{-1}$ is informative.
\end{remark}

\begin{proof}
  We compute $(p-1)! \pmod{p^2}$ by the half-factorial pairing.

  \textbf{Step 1 (Pairing).}
  Write $(p-1)! = \prod_{j=1}^m j \cdot \prod_{j=m+1}^{p-1} j$.
  Substituting $j \mapsto p-j$ in the second factor:
  $\prod_{j=m+1}^{p-1} j = \prod_{j=1}^m (p-j).$
  Hence
  \[
    (p-1)! \;=\; \prod_{j=1}^m j(p-j).
  \]

  \textbf{Step 2 (Modulo $p^2$).}
  For each $1 \leq j \leq m$, let $j^{-1}$ denote the inverse of $j$ modulo $p^2$
  (so $j \cdot j^{-1} \equiv 1 \pmod{p^2}$).  Then
  \[
    p - j = -j\bigl(1 - p j^{-1}\bigr),
  \]
  and therefore
  \[
    j(p-j) = -j^2\bigl(1 - pj^{-1}\bigr).
  \]
  Multiplying over $j = 1, \ldots, m$:
  \[
    (p-1)! = (-1)^m (m!)^2 \cdot \prod_{j=1}^m \bigl(1 - pj^{-1}\bigr).
  \]
  Since $\prod_{j=1}^m(1 - pj^{-1}) \equiv 1 - p\sum_{j=1}^m j^{-1} \pmod{p^2}$,
  we obtain
  \begin{equation}\label{eq:factorial_mod_p2}
    (p-1)! \;\equiv\; (-1)^m(m!)^2 \Bigl(1 - pS_p\Bigr) \pmod{p^2},
  \end{equation}
  where $S_p \equiv \sum_{j=1}^m j^{-1} \pmod p$ (reducing the inverse mod $p^2$
  to mod $p$).

  \textbf{Step 3 (Computing $W_p$).}
  By Wilson's theorem, $(-1)^m(m!)^2 \equiv -1 \pmod p$, so we write
  $(-1)^m(m!)^2 = -1 + p\mu_p$ for some integer $\mu_p = \bigl((-1)^m(m!)^2+1\bigr)/p$.
  Substituting into \eqref{eq:factorial_mod_p2}:
  \begin{align*}
    (p-1)! &\equiv (-1+p\mu_p)(1-pS_p)\\
           &\equiv -1 + p\mu_p + pS_p - p^2\mu_p S_p\\
           &\equiv -1 + p(\mu_p + S_p) \pmod{p^2}.
  \end{align*}
  Therefore $(p-1)! + 1 \equiv p(\mu_p + S_p) \pmod{p^2}$, and dividing by $p$:
  \[
    W_p \;=\; \frac{(p-1)!+1}{p} \;\equiv\; \mu_p + S_p \pmod p. \qedhere
  \]
\end{proof}

\begin{example}
  For $p = 5$: $m = 2$, $(m!)^2 = 4$, $(-1)^2 \cdot 4 = 4 = -1 + 5$,
  so $\mu_5 = 1$.  $S_5 = 1^{-1} + 2^{-1} = 1 + 3 = 4 \pmod 5$.
  $W_5 \equiv 1 + 4 = 5 \equiv 0 \pmod 5$. \checkmark\ ($5$ is a Wilson prime.)

  For $p = 7$: $m = 3$, $(m!)^2 = 36$, $(-1)^3 \cdot 36 = -36 = -1 + 7 \cdot (-5) = -1 - 35$,
  so $\mu_7 = -5 \equiv 2 \pmod 7$.  $S_7 = 1+4+5 = 10 \equiv 3 \pmod 7$.
  $W_7 \equiv 2+3 = 5 \pmod 7$. \checkmark\ ($W_7 = 103 \equiv 5 \pmod 7$, not a Wilson prime.)
\end{example}

% ================================================================
\section{Connection to Bernoulli Numbers}
% ================================================================

\begin{theorem}[Wilson quotient and Bernoulli numbers]\label{thm:bernoulli}
  For an odd prime $p$:
  \[
    W_p \;\equiv\; -B_{p-1} \pmod p,
  \]
  where $B_n$ denotes the $n$-th Bernoulli number.
\end{theorem}

\begin{proof}[Proof (sketch via classical results)]
  The statement $W_p \equiv -B_{p-1} \pmod p$ is a classical result
  that connects the Wilson quotient to the \emph{p-adic Bernoulli numbers}.
  We outline the two key ingredients.

  \textbf{(A) Von Staudt--Clausen and $p$-adic interpretation.}
  The von Staudt--Clausen theorem states
  $B_{p-1} + \sum_{q \text{ prime},\,(q-1)\mid(p-1)} 1/q \in \Z$.
  Since $p-1 \mid p-1$, the prime $p$ appears in the sum, so
  $p \mid \text{denom}(B_{p-1})$ and $pB_{p-1} \equiv -1 \pmod p$
  (interpreting the congruence in $\Z_p$, the ring of $p$-adic integers).
  Thus $B_{p-1}$ has a well-defined residue class modulo $p$ via
  $B_{p-1} \equiv -p^{-1} \pmod{\Z_p}$, and the statement $W_p \equiv -B_{p-1}
  \pmod p$ is understood in this $p$-adic sense.

  \textbf{(B) Connection to the half-harmonic sum.}
  From Theorem~\ref{thm:harmonic}, $W_p \equiv \mu_p + S_p \pmod p$ where
  $S_p = \sum_{j=1}^{(p-1)/2} j^{-1}$.  On the other hand, one can show
  (using the power-sum formula $\sum_{j=1}^{p-1} j^k \equiv -B_k \pmod p$
  for $(p-1) \nmid k$, together with the relation between $S_p$ and the
  power sums of inverses) that $\mu_p + S_p \equiv -B_{p-1} \pmod p$.
  A complete proof using Kummer's congruences is given in
  Ireland--Rosen \cite{Ireland1990}, Chapter~15, and in
  Berndt--Evans--Williams \cite{BEW1998}.

  \textbf{Consequence.}  A prime $p$ is a Wilson prime if and only if
  $p \mid W_p$, which by (B) is equivalent to $B_{p-1} \equiv 0 \pmod p$
  (i.e.\ $p \nmid \text{numerator}$ after reduction), or equivalently
  $p \mid pB_{p-1} + 1$.
\end{proof}

\begin{remark}
  Note that the formula $\sum_{j=1}^{p-1} j^{-1} \equiv -B_{p-2} \pmod p$
  (applying the power-sum formula with $k = p-2$, so $j^{-1} \equiv j^{p-2}$
  by Fermat) is a correct but \emph{separate} classical result.
  The connection to the Wilson quotient uses the \emph{half}-sum $S_p$
  from Theorem~\ref{thm:harmonic}, not the full sum.
\end{remark}

% ================================================================
\section{Known Wilson Primes and Heuristics}
% ================================================================

\begin{theorem}[Known Wilson primes]\label{thm:known_wilson}
  The only Wilson primes below $2 \times 10^{13}$ are $p = 5$, $p = 13$,
  and $p = 563$.
\end{theorem}

\begin{proof}
  Verification for $p = 5$:
  $(5-1)! + 1 = 24 + 1 = 25 = 5^2$, so $5^2 \mid 25$. \checkmark

  Verification for $p = 13$:
  $(13-1)! + 1 = 479001600 + 1 = 479001601 = 13^2 \cdot 2834329$,
  confirming $13^2 \mid 479001601$. \checkmark

  Verification for $p = 563$:
  $W_{563} \equiv 0 \pmod{563}$ has been verified by direct computer search.

  The non-existence of other Wilson primes below $2 \times 10^{13}$ is the
  result of an exhaustive computational search by Crandall, Dilcher, and
  Pomerance \cite{Crandall1997}, extended by Goldberg and others to $2 \times
  10^{13}$.
\end{proof}

\begin{proposition}[Heuristic density of Wilson primes]\label{prop:heuristic}
  Heuristically, the number of Wilson primes up to $x$ is approximately
  $\ln \ln x$.
\end{proposition}

\begin{proof}[Proof of heuristic]
  We model $W_p \bmod p$ as a ``random'' residue in $\{0, 1, \ldots, p-1\}$.
  The probability that a given prime $p$ is a Wilson prime is then
  approximately $1/p$.

  By the prime number theorem, the number of primes up to $x$ is $\sim x/\ln x$.
  The expected number of Wilson primes up to $x$ is therefore approximately
  \[
    \sum_{\substack{p \leq x \\ p \text{ prime}}} \frac{1}{p}
    \;\approx\; \ln \ln x,
  \]
  using Mertens' theorem $\sum_{p \leq x} 1/p = \ln \ln x + M + O(1/\ln x)$
  where $M$ is the Meissel--Mertens constant.

  This heuristic predicts that Wilson primes should be \emph{infinitely many}
  but \emph{very sparse}: approximately $\ln\ln(2 \times 10^{13}) \approx 3.4$
  Wilson primes up to $2 \times 10^{13}$, consistent with the three known ones.
\end{proof}

\begin{corollary}[Stronger congruence for Wilson primes]\label{cor:stronger}
  If $p$ is a Wilson prime, then $(p-1)! \equiv -1 \pmod{p^2}$.
\end{corollary}

\begin{proof}
  By definition, $W_p = \bigl((p-1)!+1\bigr)/p$ and $p \mid W_p$.
  This means $p^2 \mid (p-1)!+1$, i.e.\ $(p-1)! \equiv -1 \pmod{p^2}$.
\end{proof}

% ================================================================
\section{Conclusion}
% ================================================================

The Wilson quotient $W_p$ is a rich arithmetic invariant of primes.
Its connection to the harmonic sum $\sum j^{-1} \pmod p$ and to the
Bernoulli number $B_{p-1}$ places it in the framework of $p$-adic analysis
and Kummer's theory of irregular primes.

The three known Wilson primes $5$, $13$, $563$ remain isolated examples.
The heuristic prediction of $\sim \ln\ln x$ Wilson primes up to $x$ suggests
their infinite existence, but a proof remains out of reach.
Whether there exist infinitely many Wilson primes is one of the classical
open problems in analytic number theory.

\begin{thebibliography}{9}

\bibitem{Hardy1979}
G.~H.~Hardy and E.~M.~Wright.
\newblock \emph{An Introduction to the Theory of Numbers}, 5th~ed.
\newblock Oxford University Press, 1979.

\bibitem{Ireland1990}
K.~Ireland and M.~Rosen.
\newblock \emph{A Classical Introduction to Modern Number Theory}, 2nd~ed.
\newblock Springer, New York, 1990.

\bibitem{Crandall1997}
R.~Crandall, K.~Dilcher, and C.~Pomerance.
\newblock A search for Wieferich and Wilson primes.
\newblock \emph{Math.\ Comp.}, 66(217):433--449, 1997.

\bibitem{Washington1997}
L.~C.~Washington.
\newblock \emph{Introduction to Cyclotomic Fields}, 2nd~ed.
\newblock Springer, New York, 1997.

\bibitem{Berndt1994}
B.~C.~Berndt.
\newblock \emph{Ramanujan's Notebooks, Part IV}.
\newblock Springer, New York, 1994.

\bibitem{BEW1998}
B.~C.~Berndt, R.~J.~Evans, and K.~S.~Williams.
\newblock \emph{Gauss and Jacobi Sums}.
\newblock Wiley, New York, 1998.

\end{thebibliography}

\end{document}
